\subsection{Manuscript anchors for Parts B and C}

The quotations below are transcribed from the supplied page files in
the package \path{MPL_8_079-081+9+4+100}. The references use the
\emph{cote:scan} convention, followed by the manuscript page where
the transcription supplies one. Deletions have been omitted and
insertions accepted, except where an explicitly marked uncertainty
matters. These are checks of the supplied transcription, not fresh
readings of the original scans.

\paragraph{B.1: the supersingular test and the negative marginal answer.}
At 9:13 (manuscript p.12) the proposed test is:
\begin{quote}
``Cas test: soient $E_i$ $(i=1,2)$ deux courbes elliptiques sur $K$,
qui ont bonne réduction essentielle, ${}_pE_1$ isomorphe à
${}_pE_2$, $K$ contient les racines p.èmes de $1$''.
\end{quote}
The same page singles out the case where one curve reduces to
``une courbe à invariant de Hasse nul''. The next page, 9:14
(manuscript p.13), carries the marginal words ``faux [par?] Tate'';
the supplied transcription marks \emph{par} uncertain. The explicit
examples in B-RD1--2 and the ordinary theorem B-RD3 distinguish the
two cases without relying on the earlier proposed Kisin-module
construction.

\paragraph{B.2: the scope of the small-defect list.}
The hypothesis introducing the list on 9:14 (manuscript p.13) is
essential:
\begin{quote}
``Dans le cas d'une courbe elliptique, et en car. rés. $\ne 2,3$,
prenant pour $K'$ la plus petite extension sur laquelle on a
réduction semi-stable, on trouve que $G$ est un sous-groupe de
$\operatorname{Aut}(A_{s'})$, lequel est cyclique a $2,4$ ou $6$
éléments donc est premier à $p$''.
\end{quote}
This sentence does not assert the same list in residue
characteristic three. B-RD8--10 explain the additional wild
defect~12 and the exact first-level ramification ranges.

\paragraph{B.3: constant torsion and the trace normalization.}
The question at 9:14 is whether, when
``${}_pA_K$ soit un groupe constant'', the trace is
``égale à $n=\dim A_K$'' on prime-to-$p$ inertia. The rank of the
covariant Dieudonne module is $2\dim A_K$, so this is a normalization
slip; the same page correctly uses the value~$2$ for an elliptic
curve. It explicitly describes the search as
\begin{quote}
``trouver un exemple où le critère de réduction semi-stable de
Raynaud est en défaut pour $\ell=p$''.
\end{quote}
B-RD12--15 consequently distinguish the potentially good elliptic
question from the unrestricted question, where multiplicative
Tate curves are relevant.

\paragraph{B.4 and Lead B': the precise monodromy condition.}
At 9:11 (manuscript p.10) the definition is:
\begin{quote}
``satisfait à la condition de monodromie s'il admet une suite de
composition par des groupes de B.T., telle que les
$N=\operatorname{Gr}^i(M)$ aient essentiellement bonne réduction''.
\end{quote}
At 9:12 (manuscript p.11) the trace of the resulting inertia action
is called the \emph{caractère}; 9:13 asks if it is determined
``par la seule connaissance de $M(0)={}_pM$''. The comparison is
expressly restricted to elements ``d'ordre premier à $p$''.
B-RD4 shows why the filtration condition cannot silently be replaced
by potential semistability of the whole representation. B-RD17--18
prove finite upper and nonzero lower bounds for this precise
prime-to-$p$ character; the sharp general-height index remains a
separate question.

\paragraph{C.1: the asserted equivalence and the missing overlap.}
The theorem on 100:3 (manuscript p.2) assumes a finite surjective
$S'\to S$ with $S'$ geometrically unibranch and
``$\operatorname{prof\,et}_T(S)\geq2$'', and concludes:
\begin{quote}
``Alors $F$ est une équivalence des catégories''.
\end{quote}
The reduction step on 100:7 (manuscript p.6) says:
\begin{quote}
``Or $U\amalg S'\to S$ est fid. plat, et on trouve une donnée de
descente sur $A_U\amalg A'$''.
\end{quote}
The page marks this step with question marks. An actual isomorphism
on the completed overlap is an additional requirement; C-RD12
identifies it exactly after a reduced realization has been supplied.

\paragraph{C.2: nonfinite normalization was a proof obligation.}
At 100:11 (manuscript p.10) the text is explicit:
\begin{quote}
``Le cas où on ne suppose plus $S''$ fini sur $S'$ se traite par la
méthode essentiellement de $[\ ]$ p 75--76, à condition qu'on
prouve que l'isomorphisme sur $U''\simeq U'$ entre
$T_p(A'')\mid U''$ et $M_{S''}\mid U''$ se prolonge à $S''$''.
\end{quote}
The transcription records an inserted ``O.K.'' and a marginal
question mark. The missing reference is identified with
\cite[proof of Corollary~4.2, pp.75--76]{Grothendieck1966Hom}.
Its normal-base theorem does not supply the claimed reduced-base
extension. C-RD11 gives a reduced geometrically unibranch
counterexample with nonfinite normalization.

\paragraph{C.3: descent of homomorphisms of BT groups.}
At 100:25 (manuscript p.24), Corollary~3 states:
\begin{quote}
``Alors $f$ est un morphisme de descente pour la catégorie des
schémas en groupes $p$-divisibles''.
\end{quote}
Its remaining requirement is ``Il faut prouver que $X$ est
formellement net'', with $X=\underline{\operatorname{Hom}}_S(A,B)$.
C-RD1 proves the descent equalizer in the stated noetherian
setting; C-RD2 records why an unrestricted nonnoetherian rigidity
claim would be false.

\paragraph{C.4: the local-system lemma and its marginal doubt.}
At 100:17 (manuscript p.16) the proposed lemma reads:
\begin{quote}
``Un `syst local' de schémas abéliens'' sur un schéma [loc. noeth]
provient d'un schéma abélien.''
\end{quote}
The square brackets retain the transcription's uncertain expansion
of the abbreviation. The marginal judgment is unambiguous:
``Douteux, sans doute faux''. C-RD7--10 give the precise division:
effectivity over normal or geometrically unibranch bases, and under
the local finite-cover hypothesis; failure over arbitrary reduced
noetherian bases; and invariance under nilpotent reduction for an
already specified compatible local system.
